%

%

% %%%

%%%   Самарский государственный технический университет

%%%   Кафедра прикладной математики и информатики

%%%

%%%   Преамбула для статьи в журнал "Вестник Самарского государственного

%%%   технического университета. Серия: «Физико-математические науки»"

%%%   Ориентирована под MikTEx 2.4 for Win32

%%%

%%%   Хотя сам журнал собирается под GNU/Linux на дистрибутиве TeXLive



\documentclass[11pt,a4paper,twoside]{article}



\usepackage[cp1251]{inputenc}

\usepackage[T2A]{fontenc}

\usepackage[english,russian]{babel}

\usepackage{samgtubib}

\usepackage[dvips]{graphicx}

\usepackage[multiple]{footmisc}



\usepackage{samgtu}

\renewcommand{\VestnikYear}{2009} % Год издания Вестника

\renewcommand{\VestnikNumber}{2\,(19)} % Номер Вестника

\renewcommand{\VestnikMonth}{Ноябрь} % Месяц выхода Вестника

\renewcommand{\VestnikData}{01.11.09} % Дата подписания Вестника в печать

\renewcommand{\VestnikUPL}{26,64} % Усл. печ. л.

\renewcommand{\VestnikUIL}{25,73} % Уч.-изд. л.

\renewcommand{\VestnikRegNumber}{479/09} % Регистрационный номер

\renewcommand{\VestnikOrder}{994} % Номер заказа







% 

%% \usepackage{enumerate}

%%

% \begin{document}







\renewcommand{\firstpage}{26}

\renewcommand{\lastpage}{36}

% \newcommand{\bigcom}{\text{\usefont{T2A}{cmfib}{m}{n}\large{\raisebox{1pt}{,}}}} % Cимвол производной в ТФКП (ковариантная)

% \Partition{Дифференциальные уравнения}{Applied mechanics} % Раздел Вестника, в который подаётся статья

% Названия разделов см. на сайте при подаче заявки





%%%%%%%%%%%%%%%%%%%%%%%%%%%%%%%%%%%%%%%%%%%%%%%%%%%%%%%%%%%%%%%%%%%%%%%%%%%%%%%%%%%%

%Информация о статье и об авторах на русском и английском языках



% Номер УДК

\udk{517.956.326}%Обработка текста. Подготовка текстов : Языки программирования. Компьютерные языки

\msc{35L80; 35L20}% Коды MSC см. http://www.ams.org/msc/

\title{Задача со смещением для уравнения Бицадзе"--~Лыкова}% Полный заголовок

{Задача со смещением для уравнения Бицадзе"--~Лыкова}% Заголовок, который идёт в верхний колонтитул



\author{Е.\,Ю.~Арланова}% Автор(ы) для заголовка, если авторов несколько укать \inst

{А\,р\,л\,а\,н\,о\,в\,а~Е.\,Ю.}% Автор(ы) для оглавления и колонтитула через \,



\institute{\samgtu.}



\email{E-mail}{earlanova@gmail.com} %E-mail's авторов



\otherinf{\fullauthor{Екатерина Юрьевна Арланова}\regalia{к.ф.-м.н.}\position{старший преподаватель}

\dept{каф. прикладной математики и информатики}}



\keywords{уравнение Бицадзе"--~Лыкова, краевая задача, оператор Кобера"--~Эрдейи, оператор М.~Сайго, оператор Римана"--~Лиувилля, существование и единственность решения задачи}



\workabstract{Рассмотрено уравнение Бицадзе"--~Лыкова. Для этого уравнения поставлена задача со смещением с операторами Кобера"--~Эрдейи и М.~Сайго в краевом условии. Исследованы вопросы единственности (неединственности) решения задачи при различных функциях и значениях констант, входящих в краевые условия. Сформулирован и доказан ряд теорем.}



\engtitle{The problem with shift for the Bitsadze--Lykov equation }



\engauthor{E.\,Yu.~Arlanova}{A\,r\,l\,a\,n\,o\,v\,a~E.\,Yu.}





\enginstitute{\engsamgtu.}



\otherenginfm{\fullauthor{Ekaterina Yu. Arlanova} \regalia{Ph.\,D. (Phys. \& Math.)} \position{Senior Teacher} \dept{Dept. of Applied Mathematics \& Computer Science}}



\engabstract{The Bitsadze-Lykov equation is considered. The problem with shift containing the Kober--Erd\'{e}lyi and M.~Saigo operators in boundary condition is set for this equation. The questions of uniqueness (ununiqueness) of this problem solution with different functions and constants in boundary condition are investigated. The number of theorems is formulated and proved. }



\engkeywords{Bitsadze--Lykov equation, boundary value problem, Kober--Erd\'{e}lyi operator, M.~Saigo operator, Riemann--Liouville operator, existence and uniqueness of the problem solution}



\date{05/VII/2012} % Дата поступления статьи в редакцию.

\revisiondate{19/XI/2012} % В окончательном виде



%%%%%%%%%%%%%%%%%%%%%%%%%%%%%%%%%%%%%%%%%%%%%%%%%%%%%%%%%%%%%%%%%%%%%%%%%%%%%%%%%%%%





\maketitle



%Каждый пункт статьи должен начинаться с команды Название пункта

%после названия точку ставить не нужно. Пункты автоматически нумеруются.

%Если пункт нумеровать не нужно, то следует использовать в команде

% \Section необязательный параметр [n]:

%   \Section[n]{Имя пункта}

% Введение и заключение обычно не нумеруются.

Рассмотрим уравнение Бицадзе"--~Лыкова

\begin{equation}

y^2u_{xx}-u_{yy}+au_x=0, \quad |a|\leq 1,

\label{arl:1}

\end{equation}

в характеристической области 

$

D=\left\{(x,\,y):\,0<x-{y^2}/{2}<x+{y^2}/{2}<1\right\},

$

ограниченной интервалом $J=(0,\,1)$ и характеристиками данного уравнения

%\begin{multline*}

$AC=\left\{(x,\,y):\,x-{y^2}/{2}=0,\,y\leq 0\right\}$ %\\ \text{ и } 

и $BC=\left\{(x,\,y):\,x+{y^2}/{2}=1,\,y\leq 0\right\}$.

%\end{multline*}



Это уравнение, описывающее при $a>0$ процесс переноса влаги в капил\-ляр\-но-пористых средах, выведено в 1965~г. А.\,В.~Лыковым методами термодинамики необратимых процессов и с тех пор известно как уравнение влагопереноса \cite{arl:bib:13}. Ранее в 1959~г. в монографии А.\,В.~Бицадзе \cite{arl:bib:14} это уравнение приводилось в качестве примера, для которого при $|a|\leq1$ корректна по Адамару задача Коши с начальными данными на линиии $y=0$ параболического вырождения. В монографии А.\,М.~Нахушева \cite{arl:bib:15} уравнение приведено в качестве математической модели одномерного потока $u=u(x,t)$ биомассы микробной популяции и названо уравнением Бицадзе"--~Лыкова при любых значениях параметра $a\in\mathbb{R}$.



Обозначим $\Theta_0(x)=\left({x}/{2}, -\sqrt{x}\right)$ и

$\Theta_1(x)=\left(({1+x})/{2}, -\sqrt{1-x}\right)$ "--- точки пересечения

характеристик уравнения~\eqref{arl:1}, выходящих из произвольной точки

\mbox{$x\in (0, 1)$}, с характеристиками $AC$ и $BC$ соответственно. Пусть $\left(I^{\alpha}_{0+}f\right)(x)$, $\left(I^{\alpha}_{1-}f\right)(x)$ "--- левосторонний и правосторонний операторы Римана"--~Лиувилля порядка $\alpha\in\mathbb{R}$ \cite{arl:bib:10}; $(E^{\alpha, \eta}_{0+}f)(x)$ "--- левосторонний оператор Кобера"--~Эрдейи в обозначениях, принятых в монографии~\cite{arl:bib:12}; 

$(I^{\alpha, \beta, \eta}_{1-}f)(x)$ "--- правосторонний оператор М.~Сайго, введённый в~\cite{arl:bib:11}. 

Свойства этих операторов хорошо известны \cite{arl:bib:10, arl:bib:12}



Для уравнения~\eqref{arl:1} при $|a|<1$ исследуем следующую краевую задачу.



\smallskip



\begin{newthm}{Задача} Найти функцию $u(x,\,y)\in C(\overline{D})\cap C^2(D),$

удовлетворяющую уравнению~\eqref{arl:1} при $|a|<1$ в области $D$ и краевым условиям

\begin{multline}

A(x)\left(E^{\alpha_1, \frac{a-3}{4}-\alpha_1}_{0+}u[\Theta_0(t)]\right)(x)+

\\

+B(x)\left(I^{\alpha_2, \beta_2, -\alpha_2-\frac{a+3}{4}}_{1-}u[\Theta_1(t)]\right)(x)=C(x),

\label{arl:2}

\end{multline}

\begin{equation}

u(x, 0)=\tau(x),

\label{arl:3}

\end{equation}

где $A^2(x)+B^2(x)\neq 0$ для любых $x\in\overline{J};$ $A(x),$ $B(x),$ $C(x),$ $\tau(x)$ "--- заданные гладкие функции$,$ на которые в дальнейшем будут наложены определенные дополнительные условия$.$

\end{newthm}



\smallskip



Задачи, подобные сформулированной, названные позднее <<нелокальными>> или <<со смещениями>>, начиная с работ~[\citen{arl:bib:1}--\citen{arl:bib:3}]  исследовались для различных уравнений в~частных производных многими авторами (см.~обзор в~\cite{arl:bib:4}). 

В~подавляющем большинстве работ краевые условия содержали интегро"=дифференциальные операторы Римана"--~Лиувилля, например, в работе \cite{arl:bib:16}. В настоящей работе представлены некоторые результаты в направлении обобщения структуры нелокальных краевых условий.



Используя решение задачи Коши для уравнения влагопереноса при ${|a|<1}$ \cite[с. 261]{arl:bib:5}

\begin{multline*}

u(x,\,y)=\frac{\Gamma\left(\frac12\right)}{\Gamma\left(\frac{1-a}{4}\right)\Gamma\left(\frac{1+a}{4}\right)}\int^{1}_{0}\tau\left[x+\frac{y^2}{2}\left(1-2t\right)\right]\left(1-t\right)^{\frac{a-3}{4}}t^{-\frac{a+3}{4}}dt+

\\+y\frac{\Gamma\left(\frac32\right)}{\Gamma\left(\frac{3-a}{4}\right)\Gamma\left(\frac{3+a}{4}\right)}\int^{1}_{0}\nu\left[x+\frac{y^2}{2}\left(1-2t\right)\right]\left(1-t\right)^{\frac{a-1}{4}}t^{-\frac{a+1}{4}}dt,

%\label{Reshenie_Koshi}

\end{multline*}

вычислим значения $u[\Theta_0(x)]$ и $u[\Theta_1(x)]$:

$$

u[\Theta_0(x)]=k_1\left(E^{\frac{1-a}{4},\,\frac{a-3}{4}}_{0+}\tau(t)\right)(x)+k_2\sqrt{x}\left(E^{\frac{3-a}{4},\,\frac{a-1}{4}}_{0+}\nu(t)\right)(x),

$$

$$

u[\Theta_1(x)]=k_3\left(I^{\frac{1+a}{4}, 0, -\frac{a+3}{4}}_{1-}\tau(t)\right)(x)+k_4\left(I^{\frac{3-a}{4}, \frac{a-1}{4}}_{1-}\nu(t)\right)(x),

$$

где 

$$k_1=\frac{\Gamma\left(\frac12\right)}{\Gamma\left(\frac{1+a}{4}\right)},

\quad  k_2=-\frac{\Gamma\left(\frac32\right)}{\Gamma\left(\frac{3+a}{4}\right)},\quad  k_3=\frac{\Gamma\left(\frac12\right)}{\Gamma\left(\frac{1-a}{4}\right)},\quad  k_4=-\frac{\Gamma\left(\frac32\right)}{\Gamma\left(\frac{3-a}{4}\right)}.$$

Подставив $u[\Theta_0(x)]$ и $u[\Theta_1(x)]$ в краевое условие~\eqref{arl:2}, получим

\begin{multline}

A(x)k_2\left(E^{\alpha_1+\frac{3-a}{4},\frac{a-3}{4}-\alpha_1}_{0+}\sqrt{t}\nu(t)\right)(x)+

\\

+B(x)k_4\left(I^{\alpha_2+\frac{a+3}{4},\beta_2-\frac 12,-\alpha_2-\frac{a+3}{4}}_{1-}\nu(t)\right)(x)=f(x),

\label{arl:4}

\end{multline}

где 

\begin{multline*}

f(x)=C(x)-A(x)k_1\left(E^{\alpha_1+\frac{1-a}{4},\frac{a-3}{4}-\alpha_1}_{0+}\tau\right)(x)-

\\

-B(x)k_3\left(I^{\alpha_2+\frac{a+1}{4},\beta_2,-\alpha_2-\frac{a+3}{4}}_{1-}\tau\right)(x).

\end{multline*}



%Рассмотрим различные варианты изменения параметров $\alpha_1$, $\alpha_2$, $\beta_2$ и функций $A(x)$ и $B(x)$.



Справедливы следующие утверждения.



\smallskip



\begin{theorem}[1]Пусть $A(x),$ $B(x),$ $C(x)\in C^1(\overline{J});$ $\tau(x)\in C^1(\overline{J})\cap C^2(J),$ ${B(x)\neq 0}$ для любых $x\in \overline{J};$ $\alpha_2=-({a+3})/{4},$ $\beta_2=1/2,$ $\alpha_1>({a-3})/{4}.$ Тогда решение задачи \eqref{arl:1}--\eqref{arl:3} существует и единственно$.$

%\label{theorem:1}

\end{theorem}



\smallskip



\begin{proof} При выполнении условий теоремы 1 интегральное уравнение \eqref{arl:4} примет вид

\begin{equation}

B(x)k_4\nu(x)+A(x)k_2\left(E^{\alpha_1+\frac{3-a}{4},\,\frac{a-3}{4}-\alpha_1}_{0+}\sqrt{t}\nu(t)\right)(x)=f_1(x),

\label{arl:5}

\end{equation}

где

$$

f_1(x)=C(x)-A(x)k_1\left(E^{\alpha_1+\frac{1-a}{4},\,\frac{a-3}{4}-\alpha_1}_{0+}\tau\right)(x)-B(x)k_3\left(I^{-\frac 12}_{1-}\tau\right)(x).

$$

Так как $\alpha+({1-a})/{4}>0$, используя определение интегрального оператора $\left(E^{\alpha, \eta}_{0+}f\right)(x)$, запишем уравнение~\eqref{arl:5} в виде

\begin{equation}

\nu(x)+\frac{A(x)k_2}{B(x)k_4\Gamma\left(\alpha_1+\frac{3-a}{4}\right)}\int_{0}^{x}(x-t)^{\alpha_1-\frac{a+1}{4}}t^{\frac{a+5}{4}-\alpha_1}\nu(t)dt=\frac{f_1(x)}{B(x)k_4}.

\label{arl:6}

\end{equation}

Таким образом, вопрос об однозначной разрешимости задачи~\eqref{arl:1}--\eqref{arl:3} эквивалентно сводится к вопросу разрешимости уравнения~\eqref{arl:6}.



Интегральное уравнение~\eqref{arl:6} есть интегральное уравнение Вольтерры второго рода.

Для доказательства разрешимости уравнения~\eqref{arl:6} выясним гладкость правой части $f_1(x)$.



Для удобства записи обозначим

$f_1(x)=C(x)-A(x)k_1I_1(x)-B(x)k_3I_2(x),$

где $I_1(x)=\Bigl(E^{\alpha_1+\frac{1-a}{4},\,\frac{a-3}{4}-\alpha_1}_{0+}\tau\Bigr)(x)$ и $I_2(x)=\Bigl(I^{-\frac 12}_{1-}\tau(t)\Bigr)(x)$.

Рассмотрим выражение

\begin{multline*}

I_1(x)=\left(E^{\alpha_1+\frac{1-a}{4},\,\frac{a-3}{4}-\alpha_1}_{0+}\tau\right)(x)=

\\

=\frac{\sqrt{x}}{\Gamma\left(\alpha_1+\frac{1-a}{4}\right)}\int_{0}^{x}(x-t)^{\alpha_1-\frac{a+3}{4}}t^{\frac{a-3}{4}-\alpha_1}\tau(t)dt.

%\label{Ex:2:1:1:22}

\end{multline*} 

Произведя замену переменной $t=xz$, получим

$$

I_1(x)=\frac{1}{\Gamma\left(\alpha_1+\frac{1-a}{4}\right)}\int_{0}^{1}(1-z)^{\alpha_1-\frac{a+3}{4}}z^{\frac{a-3}{4}-\alpha_1}\tau(xz)dz.

$$

После несложных преобразований $I_2(x)$ можно записать так:

\begin{multline*}

I_2(x)=\frac{1}{\Gamma\left(\frac 32\right)}\int_{0}^{1}(1-x)^{\frac 12}\tau'[(1-x)t+x]\frac{dt}{\sqrt{t}}-

\\-\frac{1}{2\Gamma\left(\frac 32\right)}\int_{0}^{1}(1-x)^{-\frac 12}\tau[(1-x)t+x]\frac{dt}{\sqrt{t}}.

%\label{Ex:2:1:1:24}

\end{multline*}



Учитывая свойства функций $A(x)$, $B(x)$ и $C(x)$, явный вид $I_1(x)$ и $I_2(x)$, можно заключить, что правая часть уравнения~\eqref{arl:5} $f_1(x)\in C^2(J)$, при $x=0$ ограничена, а при $x=1$ она может обращаться в бесконечность порядка не выше~$1/2$.



Исследуем ядро $$K(x,t)=\frac{A(x)}{B(x)}(x-t)^{\alpha_1-\frac{a+1}{4}}t^{\frac{a+5}{4}-\alpha_1}.$$ На линии $t=x$ ядро имеет особенность порядка $\alpha_1-({a+1})/{4}>-1$, в остальных точках квадрата $0<x$, $t<1$  ядро дважды непрерывно дифференцируемо, так как $A(x)$, $B(x)\in C^1(\overline{J})$, $B(x)\neq 0$, $\alpha_1>({a-3})/{4}$, функция $f_1(x)\in C^2(J)$. Тогда исходя из теории интегральных уравнений Вольтерры~\cite{arl:bib:6, arl:bib:7} можно утверждать, что данное уравнение имеет единственное решение, откуда следует единственность решения данной задачи.\end{proof}



\smallskip



\begin{theorem}[2]Пусть $A(x),$ $B(x),$ $C(x)\in C^1(\overline J);$ $\tau(x)\in C^1(\overline J)\cap C^2(J);$ $A(x)\neq 0$ $\forall x\in \overline{J};$ $\alpha_1=({a-3})/{4}$, $\alpha_2>-({a+1})/{4}$, $1/2<\beta_2<1.$ Тогда решение задачи \eqref{arl:1}--\eqref{arl:3} существует и единственно$.$

%\label{theorem:2}

\end{theorem}



\smallskip



\begin{proof} При выполнении условий теоремы 2 интегральное уравнение~\eqref{arl:4} примет вид

\begin{equation}

A_1(x)k_2\nu(x)+B(x)k_4\left(I^{\alpha_2+\frac{a+3}{4},\,\beta_2-\frac 12,\,-\alpha_2-\frac{a+3}{4}}_{1-}\nu(t)\right)(x)=f_2(x),

\label{arl:7}

\end{equation}

где

$

f_2(x)=C(x)-A_1(x)k_1\Bigl(I^{-\frac 12}_{0+}\tau\Bigr)(x)-B(x)k_3\Bigl(I^{\alpha_2+\frac{a+1}{4},\,\beta_2,\,-\alpha_2-\frac{a+3}{4}}_{1-}\tau\Bigr)(x),

$

$A_1(x)=\sqrt{x}A(x)$.

Перепишем уравнение~\eqref{arl:7} следующим образом:

\begin{multline}

\nu(x)+\frac{B(x)k_4}{A_1(x)k_2\Gamma\left(\alpha_2+\frac{a+3}{4}\right)} \times \\ \times \int_{x}^{1}(t-x)^{\alpha_2+\frac{a-1}{4}}(1-t)^{-\left(\alpha_2+\beta_2+\frac{a+1}{4}\right)}\nu(t)dt=

\frac{f_2(x)}{A_1(x)k_2}.

\label{arl:8}

\end{multline}

Таким образом, вопрос об однозначной разрешимости задачи~\eqref{arl:1}--\eqref{arl:3} эквивалентно сводится к вопросу разрешимости уравнения~\eqref{arl:8}.



Интегральное уравнение~\eqref{arl:8} "--- интегральное уравнение Вольтерры второго рода.

Для доказательства разрешимости уравнения~\eqref{arl:8} выясним гладкость правой части $f_2(x)$.



Для удобства положим

$$I_3(x)=\Bigl(I^{\alpha_2+\frac{a+1}{4},\,\beta_2,\,-\alpha_2-\frac{a+3}{4}}_{1-}\tau\Bigr)(x),\quad  I_4(x)=\Bigl(I^{-\frac 12}_{0+}\tau(t)\Bigr)(x).$$ Тогда

$$f_2(x)=C(x)-B(x)k_3I_3(x)-A_1(x)k_1I_4(x).$$

Рассмотрим 

\begin{multline*}

I_3(x)=\left(I^{\alpha_2+\frac{a+1}{4},\,\beta_2,\,-\alpha_2-\frac{a+3}{4}}_{1-}\tau\right)(x)=\frac{(1-x)^{-\left(\alpha_2+\beta_2+\frac{a+1}{4}\right)}}{\Gamma\left(\alpha_2+\frac{a+1}{4}\right)}\times

\\ \times\int_{x}^{1}(t\!-\!x)^{\alpha_2\!+\!\frac{a\!-\!3}{4}}F\left(\alpha_2\!+\!\beta_2\!+\!\frac{a\!+\!1}{4},\frac{a\!+\!3}{4}\!+\!\alpha_2;\alpha_2\!+\!\frac{a\!+\!1}{4};\frac{t\!-\!x}{1\!-\!x}\right)\tau(t)dt.

%\label{Ex:2:1:1:26}

\end{multline*}

Применяя к гипергеометрической функции Гаусса формулу Больца~\mbox{\cite{arl:bib:8}}

$$F(a,\,b;\,c;\,z)=(1-z)^{c-a-b}F(c-a,\,c-b;\,c;\,z),$$

получим

\begin{multline*}

I_3(x)=\frac{\sqrt{1-x}}{\Gamma\left(\alpha_2+\frac{a+1}{4}\right)}\times

\\ \times\int_{x}^{1}\tau(t)(t\!-\!x)^{\alpha_2+\frac{a-3}{4}}(1\!-\!t)^{-\alpha_2-\beta_2-\frac{a+3}{4}}F\left(-\beta_2,-\frac 12;\alpha_2\!+\!\frac{a\!+\!1}{4};\frac{t\!-\!x}{1\!-\!x}\right)dt.

%\label{Ex:2:1:1:27}

\end{multline*}

После несложных преобразований $I_4(x)$ можно записать в виде

$$

I_4(x)=\frac{1}{2\Gamma\left(\frac 32\right)}\int_{0}^{1}\frac{\tau(xz)dz}{\sqrt{x}\sqrt{1-z}}+\frac{1}{\Gamma\left(\frac 32\right)}\int_{0}^{1}\sqrt{x}\tau'(xz)\frac{dz}{\sqrt{1-z}}.

$$



Учитывая свойства функций $A(x)$, $B(x)$ и $C(x)$, явный вид $I_3(x)$ и $I_4(x)$, можно заключить, что правая часть уравнения~\eqref{arl:7} $f_2(x)\in C^2(J)$, причём при $x=0$ она может обращаться в бесконечность порядка не выше $1/2$, а при $x=1$ "--- в бесконечность порядка не выше $\beta_2$.



Исследуем ядро $$K(x,t)=\frac{B(x)}{A(x)}(t-x)^{\alpha_2-\frac{a+1}{4}}(1-t)^{\alpha_2+\beta_2+\frac{a+1}{4}}.$$ 

На линии $t=x$ ядро имеет особенность порядка $\alpha_2-({a+1})/{4}>-1$, на линии $t=1$ "--- особенность порядка $\alpha_2+\beta_2+({a+1})/{4}>-1$, в остальных точках квадрата $0<x$, $t<1$  ядро дважды непрерывно дифференцируемо, так как $A(x)$, $B(x)\in C^1(\overline{J})$, $A(x)\neq 0$, $\alpha_2>-({a+1})/{4}$, $1/2<\beta_2<1$, функция $f(x)\hm\in C^{2}(J)$. 

Тогда исходя из теории интегральных уравнений Вольтерры~\mbox{\cite{arl:bib:6, arl:bib:7}} можно сделать вывод о том, что данное уравнение имеет единственное решение, откуда следует

единственность решения данной задачи.\end{proof}



\smallskip



\begin{theorem}[3]Пусть в условиях теоремы $2$ $\alpha_1=({a-3})/{4}$, $\alpha_2=-1/2$, $\beta_2=1/2$. Тогда решение задачи \eqref{arl:1}--\eqref{arl:3}$,$ вообще говоря$,$ не единственно$.$

\end{theorem}





\smallskip



\begin{proof} При $C(x)=0$, $\tau(x)=0$ интегральное уравнение~\eqref{arl:4} запишется в виде

\begin{equation}

\nu(x)+\frac{B(x)k_4}{A_1(x)k_2}\frac{1}{\Gamma\left(\frac{a+1}{4}\right)}\int_{x}^{1}(t-x)^{\frac{a-3}{4}}(1-t)^{-\frac{a+1}{4}}\nu(t)dt=0.

\label{arl:9}

\end{equation}

Положим $\nu_1(x)=(1-x)^{-\frac{a+1}{4}}\nu(x)$. Тогда из уравнения~\eqref{arl:9} получим интегральное уравнение Вольтерры третьего рода \cite{arl:bib:9}

\begin{equation}

(1-x)^{\frac{a+1}{4}}\nu_1(x)+\frac{B(x)k_4}{A_1(x)k_2}\frac{1}{\Gamma\left(\frac{a+1}{4}\right)}\int_{x}^{1}\frac{\nu_1(t)dt}{(t-x)^{\frac{3-a}{4}}}=0.

\label{arl:10}

\end{equation}

Пусть

$$B(x)=-\frac{A_1(x)k_2\Gamma\left(n+1\right)}{k_4\Gamma\left(n+\frac{3-a}{4}\right)}.$$

Тогда уравнение~\eqref{arl:10} примет вид

$$

(1-x)^{\frac{a+1}{4}}\nu_1(x)+a^*_1\int_{x}^{1}\frac{\nu_1(t)dt}{(t-x)^{\frac{3-a}{4}}}=0,

$$

где $$a^*_1=-\frac{\Gamma(n+1)}{\Gamma\left(\frac{a+1}{4}\right)\Gamma\left(n+\frac{3-a}{4}\right)}.$$



Покажем, что $\nu(x)=(1-x)^n$ является ненулевым решением уравнения~\eqref{arl:9}. Для этого подставляя выражение для $\nu(x)$ в уравнение~\eqref{arl:9}, делая в интеграле замену переменных $t=1-(1-x)z$, используя определение бета-функции~\cite{arl:bib:8}

$${\rm B}(x,\,y)=\int_{0}^{1}t^{x-1}(1-t)^{y-1}dt, \quad {\rm Re} x>0, \quad {\rm Re} y>0$$

и ее выражение через гамма-функции~\cite{arl:bib:8}

$${\rm B}(x,\,y)=\frac{\Gamma(x)\Gamma(y)}{\Gamma(x+y)},$$

получим верное равенство

$$(1-x)^n\left[1-\frac{\Gamma(n+1)}{\Gamma\left(\frac{a+1}{4}\right)\Gamma\left(n+\frac{3-a}{4}\right)}\cdot \frac{\Gamma\left(\frac{a+1}{4}\right)\Gamma\left(n+\frac{3-a}{4}\right)}{\Gamma(n+1)}\right]=0.$$

Отсюда следует, что $\nu(x)=(1-x)^n$ при любых $n$ удовлетворяет уравнению~\eqref{arl:9}.\end{proof}



\smallskip



\begin{theorem}[4]Пусть \mbox{$\alpha_1=({a-3})/{4},$} \mbox{$-({a+5})/{4}<\alpha_2<-({a+3})/{4},$} \mbox{$\beta_2=3/2,$} %\linebreak

$\tau(x)=(1-x)^{\delta}\tau_1(x),$ $\delta\geq\ 3/2;$ 

$A(x),$ $B(x),$ $C(x)\in C(\overline J)\cap C^1(J);$ $A(x)\neq 0,$ $B(x)\neq 0$ $\forall x\in \overline{J};$ $\tau_1(x)\in C^1(\overline J)\cap C^3(J).$ Тогда задача \eqref{arl:1}--\eqref{arl:3} имеет более одного решения$.$

%\label{theorem:4}

\end{theorem}



\begin{proof} При выполнении условий теоремы интегральное уравнение~\eqref{arl:4} примет вид

%$$

%\nu(x)+\frac{B(x)k_4}{A_1(x)k_2}\left(I^{\alpha_2+\frac{a+3}{4},\,1,\,-\alpha_2-\frac{a+3}{4}}_{1-}\nu\right)(x)=f_3(x),

%$$

%или

\begin{equation}

\nu(x)+a_{11}(x)\frac{d}{dx}\int_{x}^{1}\nu(t)(t-x)^{\alpha_2+\frac{a+3}{4}}(1-t)^{-\left(\alpha_2+\frac{a+7}{4}\right)}dt=f_3(x),

\label{arl:11}

\end{equation}

где

$$f_3(x)=\frac{f_2(x)}{A_1(x)k_2}, \quad a_{11}(x)=\frac{B(x)k_4}{A_1(x)k_2}\frac{1}{\Gamma\left(\alpha_2+\frac{a+7}{4}\right)}.$$

Покажем, что однородное уравнение, соответствующее~\eqref{arl:11}, имеет не\-три\-виальное решение.



Запишем однородное уравнение ($C(x)=0$, $\tau(x)=0$) в виде

\begin{equation}

\nu(x)+a_{11}\frac{d}{dx}\int_{x}^{1}\nu(t)(t-x)^{\alpha_2+\frac{a+3}{4}}(1-t)^{-\left(\alpha_2+\frac{a+7}{4}\right)}dt=0.

\label{arl:12}

\end{equation}

Введём новую неизвестную функцию

\begin{equation}

\varphi(x)=-\int_{1}^{x}\frac{(1-t)^{-\left(\alpha_2+\frac{a+7}{4}\right)}\nu(t)}{(t-x)^{-\left(\alpha_2+\frac{a+3}{4}\right)}}dt

\label{arl:13}

\end{equation}

и применим формулу обращения \cite[с. 39]{arl:bib:10}

$$f(x)=-\frac{\sin\pi\mu}{\pi}\frac{d}{dx}\int_{x}^{1}\frac{F(t)dt}{(t-x)^{1-\mu}}$$

интегрального уравнения Абеля

$$\int_{x}^{1}\frac{f(t)dt}{(t-x)^{\mu}}=F(t), \quad 0<\mu<1,$$

к уравнению~\eqref{arl:13}. Получим

\begin{multline}

%\begin{array}{c}

\nu(x)=-\frac{\sin\pi\left(-\left(\alpha_2+\frac{a+3}{4}\right)\right)}{\pi}(1-x)^{\alpha_2+\frac{a+7}{4}}\frac{d}{dx}\int_{x}^{1}\frac{\varphi(t)dt}{(t-x)^{\alpha_2+\frac{a+7}{4}}}=

\\

=\frac{\sin\pi\left(\alpha_2+\frac{a+3}{4}\right)}{\pi}(1-x)^{\alpha_2+\frac{a+7}{4}}\frac{d}{dx}\int_{x}^{1}\frac{\varphi(t)dt}{(t-x)^{\alpha_2+\frac{a+7}{4}}}.

%\end{array}

\label{arl:14}

\end{multline}



Найдём $\varphi(1)$. Производя в уравнении~\eqref{arl:13} замену переменных \mbox{$t{=}1{-}(1{-}x)z$}, получим

$$\varphi(x)=\int_{0}^{1}z^{-\left(\alpha_2+\frac{a+7}{4}\right)}(1-z)^{\alpha_2+\frac{a+3}{4}}\nu(1-(1-x)z)dz.$$

Тогда

\begin{multline*}

\varphi(1)=\nu(1){\rm B}\Bigl(-\alpha_2-\frac{a+3}{4},\,\alpha_2+\frac{a+7}{4}\Bigr)=

\\

=\nu(1)\frac{\Gamma\left(-\alpha_2-\frac{a+3}{4}\right)\Gamma\left(\alpha_2+\frac{a+7}{4}\right)}{\Gamma(1)}%=

%\\

=\nu(1)\frac{\pi}{\sin\pi\left(\alpha_2+\frac{a+7}{4}\right)}\neq 0 \quad (\nu(1)\neq 0).

\end{multline*}

Обозначим $\varphi(1)=C_1\neq 0$.

Если обозначить

\begin{equation}

\psi(x)=\frac{d}{dx}\varphi(x),

\label{arl:15}

\end{equation}

то очевидно

\begin{equation}

\varphi(x)=C_1-\int_{x}^{1}\psi(z)dz.

\label{arl:16}

\end{equation}



Подставляя~\eqref{arl:14} и~\eqref{arl:15} в~\eqref{arl:12}, получим

$$a_{11}(x)\psi(x)+\frac{\sin\pi\left(\alpha_2+\frac{a+3}{4}\right)}{\pi}(1-x)^{\alpha_2+\frac{a+7}{4}}\frac{d}{dx}\int_{x}^{1}\frac{\varphi(t)dt}{(t-x)^{\alpha_2+\frac{a+7}{4}}}=0.$$



Найдём $$\frac{d}{dx}\int_{x}^{1}\frac{\varphi(t)dt}{(t-x)^{\alpha_2+\frac{a+7}{4}}}.$$ 

Сделаем замену переменных $t=1-(1-x)z$ и продифференцируем:

\begin{multline*}

\frac{d}{dx}\int_{x}^{1}\frac{\varphi(t)dt}{(t-x)^{\alpha_2+\frac{a+7}{4}}}=

\frac{d}{dx}(1-x)^{-\left(\alpha_2+\frac{a+3}{4}\right)}\int_{0}^{1}\frac{\varphi(1-(1-x)z)dz}{(1-z)^{\alpha_2+\frac{a+7}{4}}}=

\\=\Bigl(\alpha_2+\frac{a+3}{4}\Bigr)(1-x)^{-\left(\alpha_2+\frac{a+7}{4}\right)}\int_{0}^{1}\frac{\varphi(1-(1-x)z)dz}{(1-z)^{\alpha_2+\frac{a+7}{4}}}+

\\+(1-x)^{-\left(\alpha_2+\frac{a+3}{4}\right)}\int_{0}^{1}\frac{\varphi'(1-(1-x)z)zdz}{(1-z)^{\alpha_2+\frac{a+7}{4}}}.

\end{multline*}

Делая замену переменных $z=({1-t})/({1-x})$ и учитывая~\eqref{arl:16}, имеем

\begin{multline*}

\frac{d}{dx}\int_{x}^{1}\frac{\varphi(t)dt}{(t-x)^{\alpha_2+\frac{a+7}{4}}}

=\Bigl(\alpha_2+\frac{a+3}{4}\Bigr)\frac{1}{1-x}\int_{x}^{1}\frac{\varphi(t)dt}{(t-x)^{\alpha_2+\frac{a+7}{4}}}+

\\

+\frac{1}{1-x}\int_{x}^{1}\frac{\varphi'(t)(1-t)dt}{(t-x)^{\alpha_2+\frac{a+7}{4}}}

=\Bigl(\alpha_2+\frac{a+3}{4}\Bigr)\frac{1}{1-x}\int_{x}^{1}\frac{C_1dt}{(t-x)^{\alpha_2+\frac{a+7}{4}}}-

\\

-\Bigl(\alpha_2+\frac{a+3}{4}\Bigr)\frac{1}{1-x}\int_{x}^{1}\frac{dt}{(t-x)^{\alpha_2+\frac{a+7}{4}}}\int_{t}^{1}\psi(z)dz

+\frac{1}{1-x}\int_{x}^{1}\frac{\psi(t)(1-t)dt}{(t-x)^{\alpha_2+\frac{a+7}{4}}}.

\end{multline*}



Рассмотрим интегралы

$$\int_{x}^{1}(t-x)^{-\left(\alpha_2+\frac{a+7}{4}\right)}dt=\left.\frac{(t-x)^{-\left(\alpha_2+\frac{a+3}{4}\right)}}{-\left(\alpha_2\frac{a+3}{4}\right)}\right|_{x}^{1}=-\frac{(1-x)^{-\left(\alpha_2+\frac{a+3}{4}\right)}}{\alpha_2+\frac{a+3}{4}},$$

\begin{multline*}

\int_{x}^{1}\frac{dt}{(t-x)^{\alpha_2+\frac{a+7}{4}}}\int_{t}^{1}\psi(z)dz

=\int_{x}^{1}\psi(z)dz\int_{x}^{z}(t-x)^{-\left(\alpha_2+\frac{a+7}{4}\right)}dt=

\\

=\int_{x}^{1}\psi(z)dz\left.\frac{(t-x)^{-\left(\alpha_2+\frac{a+3}{4}\right)}}{-\left(\alpha_2+\frac{a+3}{4}\right)}\right|_{x}^{z}

=-\frac{1}{\alpha_2+\frac{a+3}{4}}\int_{x}^{1}\psi(t)(t-x)^{-\left(\alpha_2+\frac{a+3}{4}\right)}dt.

\end{multline*}

Тогда

\begin{multline*}

\frac{d}{dx}\int_{x}^{1}\frac{\varphi(t)dt}{(t-x)^{\alpha_2+\frac{a+7}{4}}}

=-C_1(1-x)^{-\left(\alpha_2+\frac{a+7}{4}\right)}+

\\

+\frac{1}{1-x}\int_{x}^{1}\psi(t)(t-x)^{-\left(\alpha_2+\frac{a+3}{4}\right)}dt+

\frac{1}{1-x}\int_{x}^{1}\frac{\psi(t)(1-t)dt}{(t-x)^{\alpha_2+\frac{a+7}{4}}}=

\\

=-C_1(1-x)^{-\left(\alpha_2+\frac{a+7}{4}\right)}

+\int_{x}^{1}\psi(t)(t-x)^{-\left(\alpha_2+\frac{a+7}{4}\right)}dt.

\end{multline*}



Итак, получено уравнение

\begin{multline*}

a_{11}(x)\psi(x)+\frac{\sin\pi\left(\alpha_2+\frac{a+3}{4}\right)}{\pi}(1-x)^{\alpha_2+\frac{a+7}{4}}\int_{x}^{1}\psi(t)(t-x)^{-\left(\alpha_2+\frac{a+7}{4}\right)}dt=

\\=C_1\frac{\sin\pi\left(\alpha_2+\frac{a+3}{4}\right)}{\pi},

\end{multline*}

или при $B(x)\neq 0$

\begin{equation}

\psi(x)+a_{11}^*(x)\int_{x}^{1}\psi(t)(t-x)^{-\left(\alpha_2+\frac{a+7}{4}\right)}dt=h(x),

\label{arl:17}

\end{equation}

где $$a_{11}^*(x)=\frac{\sin\pi\left(\alpha_2+\frac{a+3}{4}\right)}{a_{11}(x)\pi}(1-x)^{\alpha_2+\frac{a+7}{4}},\quad  h(x)=\frac{C_1\sin\pi\left(\alpha_2+\frac{a+3}{4}\right)}{a_{11}(x)\pi}$$ $(a_{11}(x)\neq 0).$



Уравнение~\eqref{arl:17} "--- интегральное уравнение Вольтерры второго рода, для которого существует единственное решение. Отсюда следует, что исследованная задача имеет неединственное решение.



Докажем существование решения задачи.

Используя~\eqref{arl:17}, имеем уравнение

\begin{equation}

\psi(x)+a_{11}^*(x)\int_{x}^{1}\psi(t)(t-x)^{-\left(\alpha_2+\frac{a+7}{4}\right)}dt=F(x),

\label{arl:18}

\end{equation}

где $F(x)=h(x)+f_3(x)$.



Учитывая условия теоремы 4, заметим, что правая часть~\eqref{arl:18}  \mbox{$F(x)\in C(\overline{J})$}. В этом классе уравнение~\eqref{arl:18} имеет нетривиальное решение $\psi(x)$. По найденному $\psi(x)$ находится $\nu(x)$ а, следовательно, и $u(x,\,y)$.

\end{proof}



\pagebreak



%% Благодарности, гранты и прочее



\begin{thebibliography}{99}



\RBibitem{arl:bib:13}

\by Лыков~А.\,В.

\paper Применение методов термодинамики необратимых процессов с исследованием тепло и массообмена 

\jour Инж.-физ. журн.

\pages 287--304

\vol 9

\issue 3

\yr 1965

\transl англ. пер.:~

\paper Application of the methods of thermodynamics of irreversible processes to the investigation of heat and mass transfer

\by Luikov~A.\,V.

%DOI 10.1007/BF00828333

\jour J. Eng. Physics Thermophysics

\yr 1965

\vol 9

\issue 3

\pages 189--202



\RBibitem{arl:bib:14}

\by Бицадзе~А.\,В.

\book Уравнения смешанного типа 

\publ Изд-во АН СССР

\publaddr М.

\yr 1959

\totalpages 134

\transl \mbox{англ.} пер.:~

\book Equations of the Mixed Type

\by Bitsadze~A.\,V.

\publ Pergamon Press

\publaddr New York

\yr 1964

\totalpages xiii+160



\RBibitem{arl:bib:15}

\by Нахушев~А.\,М.

\book Уравнение математической биологии 

\publ Высш. шк.

\publaddr М.

\yr 1995

\totalpages 301

\misctransl

\by Nakhushev~A.\,M.

\book Equations of mathematical biology

\publaddr Moscow

\publ Vyssh. Shk.

\yr 1995

\totalpages 301



\RBibitem{arl:bib:10}

\by Самко~С.\,Г., Килбас~А.\,А., Маричев~О.\,И.

\book Интегралы и производные дробного порядка и некоторые их приложения

\publ Наука и техника

\publaddr Минск

\yr 1987

\totalpages 688

\transl англ. пер.:~

\by Samko~S.\,G., Kilbas~A.\,A., Marichev~O.\,I.

\book Fractional Integrals and Derivatives: Theory and Applications 

\publ Gordon \& Breach

\publaddr New York

\yr 1993

\totalpages xxxvi+976





\book Integrals and derivatives of fractional order and some of their applications

\publ Nauka i Tekhnika

\publaddr Minsk

\yr 1987

\totalpages 688 



\RBibitem{arl:bib:12}

\by Нахушев~А.\,М.

\book Дробное исчисление и его применение 

\publ Физматлит

\publaddr М.

\yr 2003

\totalpages 271

\misctransl

\by Nakhushev~A.\,M.

\book Fractional calculus and its applications

\publaddr Moscow

\publ Fizmatlit

\yr 2003

\totalpages 271



\Bibitem{arl:bib:11}

\by Saigo~M.

\paper A certain boundary value problem for the Euler--Darboux equation

\jour Math. Japon

\pages 377--385

\vol 24

\issue 4

\yr 1979







\RBibitem{arl:bib:1}

\by Жегалов~В.\,И.

\paper Краевая задача для уравнения смешанного типа с граничными условиями на обеих характеристиках и с~разрывами на переходной линии

\inbook Краевые задачи теории аналитических функций

\serial Уч\"eн. зап. Казан. гос. ун-та

\yr 1962

\vol 122

\issue 3

\pages 3--16

\publ Изд-во Казанского ун-та

\publaddr Казань

\mathnet{http://mi.mathnet.ru/uzku131}

\mathscinet{http://www.ams.org/mathscinet-getitem?mr=170127}

\zmath{http://www.zentralblatt-math.org/zmath/search/?an=Zbl 0161.07602}

\misctransl

\by Zhegalov~V.\,I. 

\paper A boundary-value problem for an equation of mixed type with boundary conditions on both characteristics and with discontinuities on the transition curve

\inbook Boundary value problems in the theory of analytic functions

\serial Kazan. Gos. Univ. Uchen. Zap.

\yr 1962

\vol 122

\issue 3

\pages 3--16

\publ Kazan University

\publaddr Kazan





\RBibitem{arl:bib:2}

\by Нахушев~А.\,М.

\paper Новая краевая задача для одного вырождающегося гиперболического уравнения 

\jour ДАН СССР

\pages 736--739

\vol 187

\issue 4

\yr 1969

\transl англ. пер.:~

\by Nakhushev~A.\,M.

\paper A new boundary value problem for a degenerate hyperbolic equation

\jour  Sov. Math., Dokl.

\vol 10

\pages 935--938 

\yr 1969





\RBibitem{arl:bib:3}

\by Нахушев~А.\,М.

\paper О некоторых краевых задачах для гиперболических уравнений и уравнений смешанного типа

\jour Диффер. уравн.

\pages 44--59

\vol 5

\issue 1

\yr 1969

\misctransl

\by Nakhushev~A.\,M.

\paper On Some Boundary Value Problems for Hyperbolic Equations and Equations of the Mixed Type 

\jour Differ. Uravn.

\pages 44--59

\vol 5

\issue 1

\yr 1969





\RBibitem{arl:bib:4}

\by Нахушев~А.\,М.

\book Задачи со смещением для уравнений в частных производных

\publ Наука

\publaddr М.

\yr 2006

\totalpages 287

\misctransl

\by Nakhushev~A.\,M.

\book Problems with Shift for Partial Differential Equations

\publ Nauka

\publaddr Moscow

\yr 2006

\totalpages 287



\RBibitem{arl:bib:16}

\by Репин~О.\,А., Кумыкова~С.\,К.

\paper Нелокальная задача для уравнения Бицадзе--Лыкова

\jour Изв. вузов. Матем.

\yr 2010

\issue 3

\pages 28--35

\mathnet{http://mi.mathnet.ru/ivm6709}

\mathscinet{http://www.ams.org/mathscinet-getitem?mr=2778322}

\transl англ. пер.:~

\paper A nonlocal problem for the Bitsadze-Lykov equation

\by Repin~O.\,A., Kumykova~S.\,K.

\jour Russian Math. (Iz. VUZ)

\yr 2010

\vol 54

\issue 3

\pages 24--30

\crossref{http://dx.doi.org/10.3103/S1066369X10030059}





\RBibitem{arl:bib:5}

\by Бицадзе~А.\,В.

\book Некоторые классы уравнений в частных производных

\publ Наука

\publaddr М.

\yr 1981

\totalpages 448

\misctransl

\by Bitsadze~A.\,V.

\book Some classes of partial differential equations

\publ Nauka

\publaddr Moscow

\yr 1981

\totalpages 448



\RBibitem{arl:bib:6}

\by Трикоми~Ф.

\book О~линейных уравнениях в частных производных второго порядка смешанного типа~/ пер. с итал

\publ Гостехиздат

\publaddr М.-Л.

\yr 1947

\totalpages 192

\transl англ. пер.:~

\by Tricomi~F.

\book On linear partial differential equations of the second order of mixed type

\publaddr Dayton

\publ Brown University

\yr 1948 

\totalpages 372







\RBibitem{arl:bib:7}

\by Килбас~А.\,А.

\book Интегральные уравнения: курс лекций

\publ БГУ

\publaddr Минск

\yr 2005

\totalpages 143

\misctransl

\by Kilbas~A.\,A.

\book Integral equations. A course of lectures

\publaddr Minsk

\publ BGU

\yr 2005

\totalpages 143





\Bibitem{arl:bib:8}

\by Erd\'{e}lyi~A., Magnus~W., Oberhettinger~F., Tricomi~F.\,G.

\book Higher transcendental functions

\yr 1953

\publ McGraw-Hill Book Co, Inc.

\publaddr New York -- Toronto -- London

\ed H.~Bateman

\bookvol I

\totalpages 302

\zmath{http://www.zentralblatt-math.org/zmath/search/?an=Zbl 0051.30303}

\rtransl

русск. пер.:~

\by Бейтмен~Г., Эрдейи~А.

\book Высшие трансцендентные функции. \mbox{В~3-х~т}

\yr 1973

\publ Наука

\publaddr М.

\bookvol 1

\voltitle Гипергеометрическая функция. Функция Лежандра

\totalpages 296





\RBibitem{arl:bib:9}

\by Нахушев~А.\,М.

\paper Обратные задачи для вырождающихся уравнений и интегральные уравнения Вольтерры третьего рода

\jour Диффер. уравн.

\pages 100--111

\vol 10

\issue 1

\yr 1974

\misctransl

\by Nakhushev~A.\,M.

\paper Inverse problems for degenerate equations and Volterra integral equations of the third kind

\jour  Differ. Uravn

\pages 100--111

\vol 10

\issue 1

\yr 1974





\end{thebibliography}





\noindent

\hskip6.5cm \thedate \par\noindent

\hskip6.5cm\therevdate





\makeengtitlem



\newpage



%\end{document}

